% Part: first-order-logic
% Chapter: completeness
% Section: completeness-theorem

\documentclass[../../../include/open-logic-section]{subfiles}

\begin{document}

\iftag{FOL}
      {\olfileid{fol}{com}{cth}}
      {\olfileid{pl}{com}{cth}}

\olsection{قضیهٔ تمامیت}

\begin{explain}
اکنون نتایجی را که به دست آورده‌ایم با هم ترکیب می‌کنیم و به قضیهٔ
تمامیت می‌رسیم.
\end{explain}

\begin{thm}[قضیهٔ تمامیت]
\ollabel{thm:completeness}
فرض کنید $\Gamma$ مجموعه‌ای از !!{sentence}s باشد. اگر $\Gamma$ سازگار
باشد، ارضاپذیر است.
\end{thm}

\begin{proof}
فرض کنید $\Gamma$ سازگار باشد.\iftag{FOL}{ بنا بر
  \olref[hen]{lem:henkin}، مجموعهٔ سازگار و اشباع‌شده‌ای چون
  $\Gamma' \supseteq \Gamma$ وجود دارد.}{} بنا بر \olref[lin]{lem:lindenbaum}،
مجموعه‌ای چون $\Gamma^* \supseteq \iftag{FOL}{\Gamma'}{\Gamma}$ وجود
دارد که سازگار و !!{complete} است.
\iftag{FOL}{
  چون $\Gamma' \subseteq \Gamma^*$، برای هر
  !!{formula}~$!A(x)$، مجموعهٔ $\Gamma^*$ شامل !!a{sentence} به صورت
  \iftag{prvEx}
        {$\lexists[x][!A(x)] \lif !A(c)$}
        {$\lnot\lforall[x][!A(x)] \lif \lnot !A(c)$}
  است و بنابراین $\Gamma^*$ اشباع‌شده است. اگر $\Gamma$ شامل~$\eq$
  نباشد، آنگاه بنا بر}{بنا بر}
\olref[mod]{lem:truth}،
$\iftag{FOL}{\Sat{M(\Gamma^*)}}{\pSat{v(\Gamma^*)}}{!A}$ اگر و تنها اگر
$!A \in \Gamma^*$ باشد. به‌ویژه، از این نتیجه می‌شود که برای هر
$!A \in \Gamma$، $\iftag{FOL}{\Sat{M(\Gamma^*)}}{\pSat{v(\Gamma^*)}}{!A}$؛
پس $\Gamma$ ارضاپذیر است.\iftag{FOL}{
  اگر $\Gamma$ شامل~$\eq$ باشد،
  آنگاه بنا بر \olref[ide]{lem:truth}، برای هر !!{sentence}~$!A$،
  $\Sat{\equivclass{M}{\approx}}{!A}$ اگر و تنها اگر $!A \in \Gamma^*$.
  به‌ویژه، برای هر $!A \in \Gamma$،
  $\Sat{\equivclass{M}{\approx}}{!A}$؛ پس $\Gamma$ ارضاپذیر است.}{}
\end{proof}

\begin{cor}[قضیهٔ تمامیت، صورت دوم]
\ollabel{cor:completeness}
برای هر $\Gamma$ و هر !!{sentence}~$!A$: اگر $\Gamma \Entails !A$،
آنگاه $\Gamma \Proves !A$.
\end{cor}

\begin{proof}
توجه کنید که $\Gamma$ها در \olref{cor:completeness} و
\olref{thm:completeness} به‌طور کلی سوربندی شده‌اند. برای آنکه دچار
ابهام نشویم، \olref{thm:completeness} را با متغیری دیگر بازمی‌گوییم:
برای هر مجموعهٔ $\Delta$ از !!{sentence}s، اگر $\Delta$ سازگار باشد،
ارضاپذیر است. با نقیض‌گیری، اگر $\Delta$ ارضاپذیر نباشد، $\Delta$
ناسازگار است.
از این مطلب برای اثبات نتیجه استفاده خواهیم کرد.

فرض کنید $\Gamma \Entails !A$. آنگاه بنا بر
\olref[syn][sem]{prop:entails-unsat}، مجموعهٔ
$\Gamma \cup \{\lnot !A\}$ ارضاپذیر نیست. اگر $\Gamma \cup \{\lnot
!A\}$ را به‌جای $\Delta$ بگیریم، صورت پیشین
\olref{thm:completeness} نتیجه می‌دهد که $\Gamma \cup \{\lnot !A\}$
ناسازگار است. بنا بر
\iftag{FOL}{%
  \tagrefs{prfAX/{fol:axd:prv:prop:prov-incons},
    prfSC/{fol:seq:prv:prop:prov-incons},
    prfND/{fol:ntd:prv:prop:prov-incons},
    prfTab/{fol:tab:prv:prop:prov-incons}}}{%
  \tagrefs{prfAX/{pl:axd:prv:prop:prov-incons},
    prfSC/{pl:seq:prv:prop:prov-incons},
    prfND/{pl:ntd:prv:prop:prov-incons},
    prfTab/{pl:tab:prv:prop:prov-incons}}}،
$\Gamma \Proves !A$.
\end{proof}

\tagprob{FOL}
\begin{prob}
با استفاده از \olref[fol][com][cth]{cor:completeness}،
\olref[fol][com][cth]{thm:completeness} را اثبات کنید و بدین‌ترتیب نشان
دهید که دو صورت‌بندی قضیهٔ تمامیت هم‌ارزند.
\end{prob}
\tagendprob

\tagprob{notFOL}
\begin{prob}
با استفاده از \olref[pl][com][cth]{cor:completeness}،
\olref[pl][com][cth]{thm:completeness} را اثبات کنید و بدین‌ترتیب نشان
دهید که دو صورت‌بندی قضیهٔ تمامیت هم‌ارزند.
\end{prob}
\tagendprob

\tagprob{FOL}
\begin{prob}
برای آنکه یک دستگاه !!a{derivation} کامل باشد، قواعد آن باید به‌اندازه‌ای
قوی باشند که بتوان ناسازگاری هر مجموعهٔ ارضاناپذیر را اثبات کرد. کدام‌یک
از قواعد !!{derivation} برای اثبات تمامیت لازم بودند؟ آیا هیچ‌یک از این
قواعد در هیچ جای اثبات به کار نرفته‌اند؟ برای پاسخ به این پرسش‌ها، فهرست
یا نموداری تهیه کنید که نشان دهد در هر یک از نتایج منتهی به اثبات
\olref[fol][com][cth]{thm:completeness} از کدام قواعد !!{derivation}
استفاده شده است. کاربردهای ضمنی قواعد در این اثبات‌ها را نیز حتماً ذکر کنید.
\end{prob}
\tagendprob

\tagprob{notFOL}
\begin{prob}
برای آنکه یک دستگاه !!a{derivation} کامل باشد، قواعد آن باید به‌اندازه‌ای
قوی باشند که بتوان ناسازگاری هر مجموعهٔ ارضاناپذیر را اثبات کرد. کدام‌یک
از قواعد !!{derivation} برای اثبات تمامیت لازم بودند؟ آیا هیچ‌یک از این
قواعد در هیچ جای اثبات به کار نرفته‌اند؟ برای پاسخ به این پرسش‌ها، فهرست
یا نموداری تهیه کنید که نشان دهد در هر یک از نتایج منتهی به اثبات
\olref[pl][com][cth]{thm:completeness} از کدام قواعد !!{derivation}
استفاده شده است. کاربردهای ضمنی قواعد در این اثبات‌ها را نیز حتماً ذکر کنید.
\end{prob}
\tagendprob

\end{document}
